\documentclass[11pt,a4paper]{article}

\usepackage[utf8]{inputenc}
\usepackage[T2A]{fontenc}
\usepackage[ukrainian,english]{babel}
\usepackage[top=22mm,bottom=22mm,left=22mm,right=22mm]{geometry}
\usepackage{hyperref}
\usepackage{fancyhdr}
\usepackage{parskip}
\usepackage{tcolorbox}
\usepackage{enumitem}
\usepackage{listings}
\usepackage{xcolor}

\tcbuselibrary{skins,breakable}

\hypersetup{
    pdftitle={Workshop 11 — Security \& Sandboxing: Handout},
    pdfauthor={Vadym Bondarenko},
    colorlinks=true,
    linkcolor=blue,
    urlcolor=cyan,
}

\pagestyle{fancy}
\fancyhf{}
\fancyhead[L]{\small\textit{Workshop 11 — Security \& Sandboxing}}
\fancyhead[R]{\small\thepage}
\fancyfoot[C]{\small\textit{vcdev.me · 2026}}
\renewcommand{\headrulewidth}{0.4pt}

\lstset{
    basicstyle=\ttfamily\small,
    breaklines=true,
    frame=single,
    backgroundcolor=\color{gray!10},
    rulecolor=\color{gray!40},
    xleftmargin=8pt,
    xrightmargin=8pt,
    aboveskip=6pt,
    belowskip=6pt,
}

\title{%
  {\Huge\bfseries Workshop 11: Security \& Sandboxing}\\[0.6em]
  {\large Захист по шарах --- handout}
}
\author{Vadym Bondarenko \\ \small\href{mailto:vadym.bondarenko@vcdev.me}{vadym.bondarenko@vcdev.me}}
\date{2026}

\begin{document}

\maketitle

\begin{abstract}
Пост-воркшоп референс. Defensive security: пермішени, sandbox, devcontainer, секрети, audit, supply-chain. Усі факти посилаються на офіційні docs (\href{https://code.claude.com/docs/en/security}{code.claude.com}).
\end{abstract}

\tableofcontents
\newpage

\section{Defense in depth: 3 шари}

\begin{tabular}{lll}
\hline
\textbf{Шар} & \textbf{Що робить} & \textbf{Де налаштовується} \\
\hline
Permissions & allow/deny patterns на tool-calls & \texttt{settings.json} \texttt{permissions} \\
Sandbox & fs+network ізоляція & Seatbelt/bubblewrap, \texttt{/sandbox} \\
Hooks & last-resort guards (PreToolUse blocks) & \texttt{settings.json} \texttt{hooks} \\
\hline
\end{tabular}

Композиція важлива: один шар обходять, інші ловлять.

\section{Settings.json: permissions precedence}

\begin{enumerate}
\item Enterprise (managed)
\item Personal \texttt{\textasciitilde/.claude/settings.json}
\item Project \texttt{.claude/settings.json}
\item Local override \texttt{.claude/settings.local.json}
\end{enumerate}

Deny завжди виграє над allow на тому ж рівні.

\subsection{Patterns}

\begin{lstlisting}
{
  "permissions": {
    "allow": ["Bash(git *)", "Read", "Write"],
    "deny": [
      "Bash(rm -rf *)",
      "Bash(curl * | sh)",
      "Bash(kubectl delete *)",
      "Write(.env*)"
    ]
  }
}
\end{lstlisting}

\textbf{Gotchas:}

\begin{itemize}
\item Compound commands (\texttt{rm -rf / ; whatever}) обходять прості matchers --- використай regex
\item \texttt{curl ... | sh} --- очевидний вектор; додай у deny
\item Wrappers (\texttt{sudo rm}, \texttt{xargs rm}) --- треба окрема pattern
\end{itemize}

\section{Sandboxing}

\textbf{macOS:} Seatbelt (sandbox-exec). \textbf{Linux:} bubblewrap.

\begin{lstlisting}
/sandbox enable     # одна сесія
\end{lstlisting}

Ізолює fs+network. Дозволяє Claude робити будь-що в межах sandbox без ризику для системи.

\textbf{Use case:} expermімент з ризикованим refactor, autonomous mode на чужому проекті.

\section{Devcontainer}

Reference setup: \href{https://code.claude.com/docs/en/devcontainer}{code.claude.com/docs/en/devcontainer}.

\begin{lstlisting}
// .devcontainer/devcontainer.json
{
  "image": "anthropic/claude-code:latest",
  "postCreateCommand": "./.devcontainer/init-firewall.sh",
  "remoteUser": "claude"
}
\end{lstlisting}

\texttt{init-firewall.sh} --- default-deny outbound, allowlist для Anthropic API + npm registry. Combined з \texttt{--dangerously-skip-permissions} = autonomous mode без ризику.

\section{Secrets}

\textbf{Де НЕ зберігати:}

\begin{itemize}
\item \texttt{CLAUDE.md} (комітиться в git)
\item Slide-нотатки
\item Exercise README-и
\item Hook-скрипти (стають частиною settings.json)
\end{itemize}

\textbf{Де зберігати:}

\begin{itemize}
\item \texttt{.env} (gitignore)
\item OS keychain через \texttt{apiKeyHelper}
\item Vault / OIDC для CI
\end{itemize}

\textbf{Якщо вже залилось у git:}

\begin{lstlisting}
git filter-repo --path .env --invert-paths
# rotate скомпрометований ключ
# force-push з координацією команди
\end{lstlisting}

\section{Supply-chain risks}

\begin{tabular}{lp{8cm}}
\hline
\textbf{Vector} & \textbf{Mitigation} \\
\hline
Plugin marketplace & Перевір author, source-code, recent activity, скількість installs \\
MCP server & Запускається з повними правами Claude Code; vetting обов'язковий \\
Skill з side-effects & \texttt{disable-model-invocation: true}, manual trigger \\
Hook scripts & Один \texttt{rm -rf} = катастрофа; пиши defensive shell \\
\hline
\end{tabular}

\section{Audit}

\textbf{Транскрипти:} \texttt{\textasciitilde/.claude/projects/<project-slug>/}

\begin{lstlisting}
# знайти всі bash-виклики за день
jq 'select(.tool_name == "Bash") | .input.command' \
  ~/.claude/projects/*/transcripts/*.jsonl
\end{lstlisting}

\textbf{ConfigChange hook} --- логи зміни settings.json.

\textbf{OpenTelemetry} --- structured monitoring через \href{https://code.claude.com/docs/en/monitoring-usage}{monitoring-usage}.

\section{Чек-ліст hardening}

\begin{itemize}[leftmargin=2em]
\item[$\square$] \texttt{permissions.deny} покриває \texttt{rm -rf}, deploy, db writes
\item[$\square$] \texttt{.env} у \texttt{.gitignore}, secrets не в \texttt{CLAUDE.md}
\item[$\square$] Devcontainer для autonomous mode
\item[$\square$] Plugin/MCP vetting перед install
\item[$\square$] Hook-сценарії як last-resort guards (PreToolUse + dangerous bash)
\item[$\square$] Audit-pipeline: транскрипти + OTel
\item[$\square$] API key rotation policy
\item[$\square$] OIDC / vault для headless CI
\end{itemize}

\section{Resources}

\begin{itemize}
\item \href{https://code.claude.com/docs/en/security}{code.claude.com/docs/en/security}
\item \href{https://code.claude.com/docs/en/sandboxing}{code.claude.com/docs/en/sandboxing}
\item \href{https://code.claude.com/docs/en/permissions}{code.claude.com/docs/en/permissions}
\item \href{https://code.claude.com/docs/en/devcontainer}{code.claude.com/docs/en/devcontainer}
\item \href{https://code.claude.com/docs/en/monitoring-usage}{code.claude.com/docs/en/monitoring-usage}
\item \texttt{exercises/} і \texttt{solutions/} у цій теці
\end{itemize}

\end{document}
